\chapter{数值实验与结果分析}

\section{实验设置}
采用三组实验：
\begin{itemize}
  \item \textbf{实验A：}积分精度对比（泰勒16阶 vs RK4）；
  \item \textbf{实验B：}参数估计收敛性对比（牛顿 vs 高斯牛顿）；
  \item \textbf{实验C：}噪声扰动与初值偏差鲁棒性分析。
\end{itemize}
误差指标包括均方误差（MSE）、最大绝对误差（MAE）与收敛迭代次数。

\section{积分精度结果}
\begin{table}[H]
\centering
\caption{积分方法对比结果（示例）}
\begin{tabular}{cccc}
\toprule
方法 & 步长 $h$ & 终点误差 & 运行时间（ms）\\
\midrule
Taylor-16 & $10^{-2}$ & $2.1\times 10^{-8}$ & 18.4 \\
RK4 & $10^{-2}$ & $4.6\times 10^{-6}$ & 15.2 \\
Taylor-16 & $5\times 10^{-3}$ & $3.7\times 10^{-10}$ & 35.7 \\
RK4 & $5\times 10^{-3}$ & $2.8\times 10^{-7}$ & 30.1 \\
\bottomrule
\end{tabular}
\end{table}
结果显示在同等步长下，高阶泰勒法可显著降低截断误差。

\section{参数估计结果}
\begin{table}[H]
\centering
\caption{牛顿法与高斯牛顿法对比（示例）}
\begin{tabular}{ccccc}
\toprule
方法 & 初值偏差 & 迭代次数 & 最终目标值 & 是否收敛\\
\midrule
Newton & 小 & 7 & $1.3\times10^{-6}$ & 是 \\
Gauss--Newton & 小 & 9 & $1.5\times10^{-6}$ & 是 \\
Newton & 中 & 20 & $8.2\times10^{-4}$ & 否（震荡） \\
Gauss--Newton(阻尼) & 中 & 15 & $2.4\times10^{-5}$ & 是 \\
\bottomrule
\end{tabular}
\end{table}

\section{鲁棒性讨论}
综合实验可得：
\begin{enumerate}
  \item 噪声水平升高时，牛顿法对海森矩阵条件数更敏感；
  \item 阻尼高斯牛顿法在“初值偏差中等 + 中等噪声”场景更稳健；
  \item 采用自适应步长可在保证精度的同时控制总体算时。
\end{enumerate}
